% !PRINTED 332
\weberpage{332}

%<<CONT
\weq{2}{\begin{aligned}
\Delta'_\xi &= \Delta_{\xi+1} - \Delta_\xi. \\
&\cdot\ \cdot\ \cdot\ \cdot\ \cdot\ \cdot\ \cdot\ \cdot\ \cdot\ \cdot \\
\Delta_\xi^{(n-1)} &= \Delta_{\xi+1}^{(n-2)} - \Delta_\xi^{(n-2)}
\end{aligned}}

is set,

\weq{3}{\varphi(\xi) = \varphi(0) + \Delta_0 B_1^{(\xi)} + \Delta'_0 B_2^{(\xi)} + \cdots + \Delta_0^{(n-1)} B_n^{(\xi)}.}

The value of the formula lies in this, that the sequence of differences $\Delta_0, \Delta'_0, \Delta''_0 \ldots$ in many cases presents numerical values decreasing so rapidly that one can content oneself with a few of the first terms of the series (3).

Now if $f(x)$ is the function to be investigated and $(\alpha, \beta)$ the interval in which one of the sought roots lies, then, understanding by $m$ any suitable integer, we set
\[
\delta = \frac{\beta - \alpha}{m}; \qquad x = \alpha + \delta \xi,
\]
and if we accordingly set
\[
\begin{aligned}
\Delta_x &= f(x + \delta) - f(x) \\
\Delta'_x &= \Delta_{x+\delta} - \Delta_x \\
&\cdot\ \cdot\ \cdot\ \cdot\ \cdot\ \cdot\ \cdot\ \cdot
\end{aligned}
\]
then

\weq{4}{f(x) = f(\alpha) + \frac{\Delta_\alpha}{\delta}(x - \alpha) + \frac{\Delta'_\alpha}{\delta^2}\frac{(x-\alpha)(x-\alpha-\delta)}{2} + \cdots}

One could just as well start from the other end point $\beta$ of the interval, and would then have to set in (3)
\[
x = \beta - \delta \xi
\]
One would then obtain
\[
\begin{aligned}
\Delta_x &= f(x - \delta) - f(x) \\
\Delta'_x &= \Delta_{x-\delta} - \Delta_x \\
&\cdot\ \cdot\ \cdot\ \cdot\ \cdot\ \cdot\ \cdot\ \cdot
\end{aligned}
\]

\weq{5}{f(x) = f(\beta) + \frac{\Delta_\beta}{\delta}(x - \beta) + \frac{\Delta'_\beta}{\delta^2}\frac{(x-\beta)(x-\beta+\delta)}{2} + \cdots}

If one interprets the equation
\[
y = f(x)
\]
as the equation of a curve in a rectangular coordinate system $x, y$, then the relations can be illustrated geometrically.

If, for example, one applies equation (4) to the interval from $\alpha$ to $(\alpha + \delta)$ and contents oneself with taking account of the first difference, that means, in the language of geometry,
%CONT>>
